\section{Jailbreak Defenses}
\label{sec:jailbreak-defenses}

Chương trước đã trình bày các nhóm jailbreak attacks và cho thấy rằng
attack có thể được xây dựng thủ công, tự động hóa, tối ưu bằng thuật toán
hoặc thích nghi với cơ chế phòng thủ. Chương này phân tích các nhóm
phòng thủ trước jailbreak ở cấp đầu vào, cấp mô hình, thời điểm suy luận,
đầu ra và toàn bộ kiến trúc hệ thống.

Mục tiêu của jailbreak defense không chỉ là giảm tỷ lệ mô hình tạo phản
hồi không an toàn, mà còn phải duy trì khả năng hỗ trợ các yêu cầu hợp lệ.
Do đó, một defense cần được đánh giá đồng thời theo safety, utility,
over-refusal, chi phí suy luận và khả năng chống lại adaptive attacks.

\subsection{Taxonomy các phương pháp phòng thủ}
\label{subsec:defense-taxonomy}

Có thể phân loại jailbreak defenses theo vị trí triển khai trong pipeline
của hệ thống LLM. Bảng~\ref{tab:jailbreak-defense-taxonomy} tóm tắt các
nhóm phòng thủ chính được phân tích trong chương này.

\begin{table}[H]
\centering
\caption{Taxonomy các phương pháp phòng thủ trước jailbreak.}
\label{tab:jailbreak-defense-taxonomy}
\begin{tabularx}{\textwidth}{
    >{\bfseries}p{0.24\textwidth}
    p{0.30\textwidth}
    X
}
\toprule
Nhóm phòng thủ
& Vị trí triển khai
& Mục tiêu chính \\
\midrule

Input filtering
& Trước khi prompt được đưa vào mô hình
& Phát hiện hoặc chặn đầu vào có rủi ro, bao gồm prompt bất thường hoặc
có dấu hiệu jailbreak. \\

Prompt preprocessing
& Trước hoặc trong quá trình xây dựng context
& Chuẩn hóa, làm sạch hoặc tái diễn đạt đầu vào để giảm tác động của
obfuscation và prompt injection. \\

Adversarial training
& Trong quá trình huấn luyện hoặc fine-tuning
& Bổ sung dữ liệu tấn công vào huấn luyện để tăng độ bền của hành vi từ
chối. \\

Inference-time defenses
& Trong quá trình suy luận
& Kiểm tra tính ổn định của phản hồi, lấy mẫu nhiều lần hoặc dùng cơ chế
phát hiện rủi ro tại runtime. \\

Output moderation
& Sau khi mô hình sinh phản hồi
& Kiểm duyệt phản hồi trước khi gửi cho người dùng. \\

Guard models
& Xung quanh mô hình chính
& Dùng một mô hình hoặc classifier riêng để đánh giá input, output hoặc
trajectory. \\

System-level defenses
& Cấp ứng dụng và kiến trúc
& Giới hạn quyền công cụ, sandbox, xác nhận người dùng và monitoring. \\

Defense-in-depth
& Toàn bộ pipeline
& Kết hợp nhiều lớp phòng thủ để tránh phụ thuộc vào một cơ chế duy nhất. \\

\bottomrule
\end{tabularx}
\end{table}

Cách phân loại này nhấn mạnh rằng defense không nhất thiết nằm bên trong
LLM. Trong các hệ thống thực tế, độ an toàn thường phụ thuộc vào nhiều
thành phần phối hợp: policy, guard model, system prompt, mô hình chính,
bộ điều khiển công cụ, kiểm duyệt đầu ra và giám sát sau triển khai.

\subsection{Input filtering}
\label{subsec:input-filtering}

Input filtering là lớp phòng thủ được đặt trước mô hình chính. Mục tiêu
là phát hiện hoặc chặn các prompt có dấu hiệu rủi ro trước khi chúng được
đưa vào LLM. Các bộ lọc đầu vào có thể dựa trên luật, từ khóa, classifier,
embedding similarity hoặc một guard model chuyên đánh giá an toàn.

Một input filter đơn giản có thể kiểm tra xem yêu cầu có thuộc nhóm chính
sách không được phép hay không. Các hệ thống phức tạp hơn có thể đánh giá
cả ý định ngữ nghĩa, mức độ bất thường của prompt, dấu hiệu obfuscation
hoặc các chỉ dẫn cố gắng thay đổi vai trò của mô hình.

Ưu điểm của input filtering là có thể ngăn một số attack ngay từ đầu, giảm
tải cho mô hình chính và dễ tích hợp vào pipeline triển khai. Tuy nhiên,
input filtering có ba hạn chế quan trọng. Thứ nhất, nếu bộ lọc dựa quá
nhiều vào đặc trưng bề mặt, attacker có thể thay đổi cách diễn đạt để né
lọc. Thứ hai, bộ lọc quá nghiêm ngặt có thể làm tăng over-refusal. Thứ ba,
input filtering khó xử lý các trường hợp multi-turn hoặc indirect prompt
injection, nơi rủi ro chỉ xuất hiện khi kết hợp nhiều phần ngữ cảnh.

Do đó, input filtering nên được xem là một lớp phòng thủ ban đầu, không
phải cơ chế bảo vệ duy nhất.

\subsection{Prompt preprocessing}
\label{subsec:prompt-preprocessing}

Prompt preprocessing là nhóm kỹ thuật xử lý đầu vào trước khi đưa vào mô
hình. Mục tiêu là giảm tác động của những biến đổi bề mặt như obfuscation,
encoding hoặc cách trình bày gây nhiễu.

Các thao tác preprocessing có thể bao gồm chuẩn hóa ký tự, loại bỏ nhiễu
không cần thiết, tách rõ phần instruction và phần data, hoặc chuyển prompt
về một định dạng có cấu trúc hơn. Với các hệ thống retrieval hoặc
tool-integrated agents, preprocessing còn có thể đánh dấu nội dung bên
ngoài là dữ liệu không đáng tin cậy, thay vì cho phép nó được hiểu như chỉ
dẫn điều khiển hệ thống.

Một nguyên tắc quan trọng là không trộn lẫn dữ liệu bên ngoài với chỉ dẫn
cấp hệ thống. Ví dụ, nội dung từ website, tài liệu, email hoặc kết quả trả
về từ công cụ nên được đặt trong vùng ngữ cảnh được đánh dấu rõ là
external content. Cách này không loại bỏ hoàn toàn prompt injection, nhưng
giúp mô hình và các guard dễ phân biệt hơn giữa instruction và data.

Hạn chế của prompt preprocessing là nó không thể bảo đảm giữ nguyên toàn
bộ ý nghĩa của prompt. Nếu preprocessing làm thay đổi nội dung quá mạnh,
hệ thống có thể hiểu sai yêu cầu hợp lệ. Ngược lại, nếu preprocessing quá
nhẹ, nó có thể không đủ để loại bỏ những biến đổi adversarial.

\subsection{Adversarial training}
\label{subsec:adversarial-training}

Adversarial training là cách đưa các ví dụ tấn công vào quá trình huấn
luyện hoặc fine-tuning để tăng độ bền của mô hình. Thay vì chỉ huấn luyện
mô hình trên các yêu cầu thông thường, hệ thống được tiếp xúc thêm với
những biến thể jailbreak và học cách phản hồi an toàn.

Một tập dữ liệu adversarial training có thể được biểu diễn dưới dạng

\begin{equation}
    \mathcal{D}_{\mathrm{adv}}
    =
    \left\{
        (x_{\mathrm{adv}}^{(i)}, y_{\mathrm{safe}}^{(i)})
    \right\}_{i=1}^{N},
    \label{eq:adversarial-training-dataset}
\end{equation}

trong đó $x_{\mathrm{adv}}^{(i)}$ là đầu vào có tính adversarial và
$y_{\mathrm{safe}}^{(i)}$ là phản hồi an toàn mong muốn. Mô hình được tối
ưu để duy trì hành vi từ chối hoặc chuyển hướng phù hợp trước các biến thể
đầu vào này.

Ưu điểm của adversarial training là nó cải thiện trực tiếp hành vi của mô
hình, thay vì chỉ đặt thêm một lớp lọc bên ngoài. Nếu tập attack đủ đa
dạng, mô hình có thể học được các đặc trưng ngữ nghĩa của yêu cầu không an
toàn thay vì chỉ nhận diện mẫu bề mặt.

Tuy nhiên, adversarial training có hai thách thức lớn. Thứ nhất, không
gian jailbreak gần như không giới hạn, nên dữ liệu huấn luyện khó bao phủ
được mọi biến thể. Thứ hai, nếu huấn luyện quá thiên về từ chối, mô hình
có thể tăng over-refusal và làm giảm utility trên yêu cầu hợp lệ. Vì vậy,
adversarial training cần đi kèm đánh giá trên cả harmful prompts và benign
prompts.

\subsection{Inference-time defenses}
\label{subsec:inference-time-defenses}

Inference-time defenses là các phương pháp được áp dụng trong quá trình
suy luận, không nhất thiết phải thay đổi tham số của mô hình. Nhóm này
đặc biệt hữu ích khi nhà triển khai không có khả năng fine-tune mô hình
hoặc muốn bổ sung lớp bảo vệ độc lập với model provider.

Một hướng tiêu biểu là kiểm tra độ ổn định của mô hình trước các biến thể
nhỏ của prompt. SmoothLLM dựa trên quan sát rằng một số adversarial prompts
có tính giòn trước các thay đổi ký tự nhỏ. Phương pháp này tạo nhiều bản
biến đổi của cùng một prompt, truy vấn mô hình trên các bản đó, rồi tổng
hợp kết quả để phát hiện đầu vào đáng ngờ \cite{robey2023smoothllm}.

Ở mức khái quát, một inference-time defense dạng smoothing có thể được mô
tả như sau. Với prompt đầu vào $x$, hệ thống tạo ra $K$ biến thể:

\begin{equation}
    \tilde{x}_1, \tilde{x}_2, \ldots, \tilde{x}_K
    \sim
    \mathcal{P}(x),
    \label{eq:prompt-perturbation}
\end{equation}

trong đó $\mathcal{P}$ là phân phối biến đổi. Sau đó hệ thống đánh giá các
phản hồi tương ứng và đưa ra quyết định cuối cùng bằng hàm tổng hợp:

\begin{equation}
    d
    =
    g
    \left(
        F_{\theta}(\tilde{x}_1),
        F_{\theta}(\tilde{x}_2),
        \ldots,
        F_{\theta}(\tilde{x}_K)
    \right).
    \label{eq:smoothing-aggregation}
\end{equation}

Các biến thể sau này như SemanticSmooth mở rộng ý tưởng smoothing từ biến
đổi ký tự sang biến đổi ngữ nghĩa, nhằm tăng độ bền trước các attack không
chỉ phụ thuộc vào token-level perturbation \cite{ji2024semanticsmooth}.

Ưu điểm của inference-time defenses là có thể triển khai mà không cần huấn
luyện lại mô hình chính. Tuy nhiên, chúng thường làm tăng chi phí suy luận
vì cần nhiều lần gọi mô hình hoặc guard model. Ngoài ra, nếu attacker biết
cơ chế phòng thủ, họ có thể thiết kế adaptive attack để tạo đầu vào ổn định
hơn qua các biến thể.

\subsection{Output moderation}
\label{subsec:output-moderation}

Output moderation là lớp kiểm tra phản hồi sau khi mô hình đã sinh nội
dung nhưng trước khi gửi cho người dùng. Lớp này có thể dùng rule-based
filter, classifier hoặc guard model để đánh giá xem phản hồi có vi phạm
chính sách hay không.

Output moderation có vai trò quan trọng vì input filtering không phải lúc
nào cũng dự đoán được phản hồi cuối cùng. Một prompt ban đầu có thể trông
hợp lệ nhưng mô hình lại sinh nội dung không phù hợp; ngược lại, một prompt
có dấu hiệu rủi ro có thể vẫn dẫn đến phản hồi an toàn. Do đó, kiểm tra
đầu ra giúp giảm rủi ro ở bước cuối của pipeline.

Tuy nhiên, output moderation cũng có hạn chế. Nếu phản hồi đã được sinh
đầy đủ, hệ thống đã tiêu tốn chi phí suy luận trước khi phát hiện lỗi.
Ngoài ra, nếu output guard quá nghiêm ngặt, nó có thể chặn các phản hồi hợp
lệ, đặc biệt trong các ngữ cảnh giáo dục, nghiên cứu, y tế hoặc an toàn
công cộng. Vì vậy, output moderation cần phân biệt giữa nội dung gây hại
và nội dung phân tích an toàn ở mức phù hợp.

Trong agentic systems, output moderation không nên chỉ kiểm tra văn bản
cuối cùng. Hệ thống còn cần kiểm tra các lời gọi công cụ, tham số hành động
và trạng thái trung gian, vì rủi ro có thể xảy ra trước khi phản hồi cuối
được tạo ra.

\subsection{Representation-level defenses}
\label{subsec:representation-level-defenses}

Representation-level defenses cố gắng phát hiện hoặc can thiệp vào các
biểu diễn bên trong của mô hình, thay vì chỉ kiểm tra prompt hoặc output ở
dạng văn bản. Ý tưởng chung là các yêu cầu có rủi ro hoặc adversarial có
thể tạo ra những mẫu kích hoạt khác biệt trong hidden states, attention
patterns hoặc embedding space.

Một hướng tiếp cận là huấn luyện classifier trên biểu diễn trung gian của
mô hình để phân biệt giữa benign prompts và adversarial prompts. Một hướng
khác là can thiệp vào activation để làm giảm khả năng mô hình đi theo
hành vi không mong muốn.

Ưu điểm của representation-level defenses là chúng có thể phát hiện những
tín hiệu không dễ thấy ở bề mặt văn bản. Tuy nhiên, nhóm phương pháp này
thường yêu cầu quyền truy cập sâu vào mô hình, khó áp dụng với API đóng và
có thể phụ thuộc mạnh vào kiến trúc cụ thể. Ngoài ra, việc can thiệp vào
biểu diễn bên trong có thể ảnh hưởng đến utility nếu không được hiệu chỉnh
cẩn thận.

Vì vậy, representation-level defenses phù hợp hơn với bối cảnh nghiên cứu
hoặc mô hình mã nguồn mở, trong khi các hệ thống thương mại đóng thường
dựa nhiều hơn vào guard model, input/output moderation và system-level
controls.

\subsection{Guard models}
\label{subsec:guard-models}

Guard model là một mô hình phụ được dùng để đánh giá mức độ an toàn của
input, output hoặc hành động trung gian. Khác với rule-based filter, guard
model có thể học biểu diễn ngữ nghĩa phức tạp hơn và xử lý những prompt
không khớp với mẫu cố định.

Có thể triển khai guard model ở nhiều vị trí:

\begin{enumerate}[leftmargin=1.2cm]
    \item \textbf{Input guard:}
    đánh giá prompt trước khi đưa vào LLM chính;

    \item \textbf{Output guard:}
    đánh giá phản hồi trước khi gửi cho người dùng;

    \item \textbf{Tool guard:}
    kiểm tra lời gọi công cụ và tham số hành động;

    \item \textbf{Trajectory guard:}
    đánh giá toàn bộ chuỗi quan sát, kế hoạch và hành động của agent.
\end{enumerate}

Guard models giúp tách nhiệm vụ sinh nội dung khỏi nhiệm vụ đánh giá an
toàn. Điều này có lợi vì mô hình chính có thể tối ưu cho khả năng hỗ trợ
người dùng, trong khi guard model tập trung vào chính sách an toàn. Tuy
nhiên, guard model cũng có thể bị lỗi, bị né tránh hoặc tạo ra over-refusal.

Một vấn đề quan trọng là guard model cần được đánh giá độc lập. Nếu guard
được dùng làm evaluator trong benchmark, nhưng chính guard đó cũng được
dùng trong hệ thống phòng thủ, kết quả đánh giá có thể bị thiên lệch. Do
đó, các benchmark như HarmBench nhấn mạnh nhu cầu chuẩn hóa dữ liệu, attack
artifacts và phương pháp chấm điểm khi so sánh các defense
\cite{mazeika2024harmbench}.

\subsection{Defense-in-depth}
\label{subsec:defense-in-depth}

Không có một cơ chế phòng thủ đơn lẻ nào đủ để bảo vệ hoàn toàn hệ thống
LLM trước jailbreak. Input filtering có thể bỏ sót các prompt được diễn
đạt tinh vi. Safety alignment có thể không bao phủ mọi biến thể đầu vào.
Output moderation có thể phát hiện muộn. Guard models có thể sai hoặc bị
attacker thích nghi. Vì vậy, hệ thống thực tế thường cần defense-in-depth,
tức là kết hợp nhiều lớp phòng thủ.

Hình~\ref{fig:layered-safety-architecture-defense} minh họa một kiến trúc
phòng thủ nhiều lớp cho hệ thống LLM có khả năng gọi công cụ.

\begin{figure}[H]
\centering
\begin{tikzpicture}[
    >=Latex,
    node distance=0.65cm,
    component/.style={
        draw,
        rounded corners=3pt,
        align=center,
        minimum width=2.35cm,
        minimum height=1.0cm,
        text width=2.1cm,
        inner sep=5pt
    },
    external/.style={
        draw,
        dashed,
        rounded corners=3pt,
        align=center,
        minimum width=2.35cm,
        minimum height=1.0cm,
        text width=2.1cm,
        inner sep=5pt
    },
    arrow/.style={
        ->,
        thick
    }
]

\node[external] (user) {
    \textbf{Người dùng}\\
    Prompt đầu vào
};

\node[component, right=of user] (inputguard) {
    \textbf{Input guard}\\
    Kiểm tra prompt
};

\node[component, right=of inputguard] (llm) {
    \textbf{Aligned LLM}\\
    Sinh phản hồi
};

\node[component, right=of llm] (outputguard) {
    \textbf{Output guard}\\
    Kiểm tra output
};

\node[external, right=of outputguard] (response) {
    \textbf{Phản hồi}\\
    Đến người dùng
};

\node[component, below=1.05cm of llm] (controller) {
    \textbf{Action controller}\\
    Kiểm tra công cụ
};

\node[external, below=1.05cm of controller] (tools) {
    \textbf{Tools/Environment}\\
    Dữ liệu và hành động
};

\node[component, below=1.05cm of inputguard] (monitor) {
    \textbf{Monitoring}\\
    Log và cảnh báo
};

\draw[arrow] (user) -- (inputguard);
\draw[arrow] (inputguard) -- (llm);
\draw[arrow] (llm) -- (outputguard);
\draw[arrow] (outputguard) -- (response);

\draw[arrow] (llm) -- (controller);
\draw[arrow] (controller) -- (tools);
\draw[arrow, bend right=25] (tools.east) to (llm.south east);

\draw[arrow, dashed] (inputguard) -- (monitor);
\draw[arrow, dashed] (controller.west) to (monitor.east);

\end{tikzpicture}
\caption{Kiến trúc phòng thủ nhiều lớp cho hệ thống LLM có khả năng gọi
công cụ.}
\label{fig:layered-safety-architecture-defense}
\end{figure}

Trong kiến trúc này, input guard giảm rủi ro trước khi prompt đến mô hình.
Aligned LLM chịu trách nhiệm sinh phản hồi theo chính sách đã học.
Output guard kiểm tra nội dung trước khi trả về người dùng. Action
controller giới hạn quyền của công cụ, kiểm tra tham số hành động và có
thể yêu cầu xác nhận cho các thao tác quan trọng. Monitoring ghi nhận log,
phát hiện bất thường và hỗ trợ phân tích sau sự cố.

Defense-in-depth đặc biệt quan trọng đối với agentic systems. Khi LLM có
thể gọi công cụ, truy cập dữ liệu hoặc thực hiện hành động, rủi ro không
chỉ nằm ở văn bản cuối cùng. Một hệ thống an toàn cần kiểm soát cả prompt,
context, kế hoạch, lời gọi công cụ, output và trạng thái môi trường.

\subsection{Hạn chế của các phương pháp phòng thủ}
\label{subsec:defense-limitations}

Mặc dù các phương pháp phòng thủ có thể giảm rủi ro jailbreak, chúng vẫn
có những hạn chế chung.

Thứ nhất, không gian đầu vào ngôn ngữ tự nhiên rất lớn. Một defense được
đánh giá tốt trên một tập attack cố định có thể không bền vững trước các
biến thể mới hoặc adaptive attacks. Do đó, kết quả trên benchmark cần được
hiểu như bằng chứng thực nghiệm, không phải bảo đảm tuyệt đối.

Thứ hai, safety và utility thường có sự đánh đổi. Một defense quá nghiêm
ngặt có thể làm giảm unsafe compliance nhưng đồng thời tăng over-refusal.
Điều này làm hệ thống kém hữu ích, đặc biệt trong các ngữ cảnh mà người
dùng có nhu cầu hợp lệ như học tập, nghiên cứu hoặc phân tích rủi ro.

Thứ ba, nhiều defense làm tăng chi phí triển khai. Inference-time defenses
có thể cần nhiều lần gọi mô hình; guard models cần tài nguyên riêng;
monitoring và human confirmation làm tăng độ phức tạp vận hành. Vì vậy,
lựa chọn defense cần cân bằng giữa mức rủi ro, chi phí và yêu cầu ứng dụng.

Thứ tư, defense có thể bị attacker thích nghi. Nếu attacker biết hệ thống
dùng input guard, output guard hoặc smoothing, họ có thể điều chỉnh prompt
để tối ưu qua các lớp đó. Vì vậy, một defense đáng tin cậy cần được đánh
giá trong cả thiết lập static và adaptive.

\subsection{Tổng kết chương}
\label{subsec:jailbreak-defenses-summary}

Chương này đã trình bày các nhóm phòng thủ chính trước jailbreak, bao gồm
input filtering, prompt preprocessing, adversarial training,
inference-time defenses, output moderation, representation-level defenses,
guard models và defense-in-depth. Mỗi nhóm defense có ưu điểm và hạn chế
riêng, phụ thuộc vào vị trí triển khai, quyền truy cập mô hình, chi phí
suy luận và loại attack cần chống lại.

Một kết luận quan trọng là không nên phụ thuộc vào một lớp phòng thủ duy
nhất. Jailbreak là một bài toán động, trong đó attacker có thể thay đổi
cách diễn đạt, tự động hóa tìm kiếm prompt hoặc thích nghi với defense.
Do đó, các hệ thống thực tế cần kết hợp nhiều lớp bảo vệ và đánh giá chúng
trên cả safety, utility, over-refusal, chi phí và khả năng chống adaptive
attacks.

Chương tiếp theo trình bày các phương pháp evaluation và benchmark dùng để
đo lường attack success, refusal, harmfulness, utility và độ bền vững của
các defense.